%=============================================================================
% WAVE 4, ITERATION 16: DEFINITIVE EXPERIMENTAL PROPOSALS
% Actionable documents for experimentalists. Full specifications.
%
% Agent 17 (Experimental Proposals) | DFD Research Programme | Model: Opus 4.6
% Building on: W4i15_Experimental_update.tex, W4i14_Experimental_update.tex,
%              W4i10_P11_SQMS_Proposal.tex, W4i13_Experimental_update.tex,
%              W4i15_P11_update.tex, ITERATION_TRACKER.md (all iterations)
%
% Status flags: [VERIFIED], [NEW], [ACTIONABLE], [CORRECTED], [CRITICAL]
% Date: 20 March 2026
%=============================================================================

\documentclass[11pt]{article}
\usepackage[utf8]{inputenc}
\usepackage{amsmath,amssymb,amsfonts,mathtools}
\usepackage{geometry}
\geometry{margin=1in}
\usepackage{booktabs}
\usepackage{siunitx}
\usepackage{hyperref}
\usepackage{xcolor}
\usepackage{enumitem}
\usepackage{float}
\usepackage{array}
\usepackage{longtable}
\usepackage{tcolorbox}
\tcbuselibrary{skins,breakable}

\newcommand{\ee}[1]{\times 10^{#1}}
\newcommand{\lbone}{|\lambda - 1|}
\newcommand{\dpsi}{\delta\psi}
\newcommand{\Qpsi}{Q_\psi}
\newcommand{\Qmech}{Q_{\rm mech}}
\newcommand{\SNR}{\mathrm{SNR}}
\newcommand{\LCDM}{\Lambda\text{CDM}}
\newcommand{\Geff}{G_{\rm eff}}

\definecolor{proposalcolor}{rgb}{0.0,0.35,0.55}
\definecolor{actioncolor}{rgb}{0.0,0.5,0.0}
\definecolor{newcolor}{rgb}{0.6,0.0,0.6}
\definecolor{warncolor}{rgb}{0.8,0.2,0.0}

\newcommand{\confirmed}[1]{\textcolor{blue}{\textbf{CONFIRMED:} #1}}
\newcommand{\corrected}[1]{\textcolor{red}{\textbf{CORRECTED:} #1}}
\newcommand{\newfinding}[1]{\textcolor{teal}{\textbf{NEW FINDING:} #1}}
\newcommand{\caution}[1]{\textcolor{orange}{\textbf{CAUTION:} #1}}

\title{Wave 4, Iteration 16: Definitive Experimental Proposals\\
\large Actionable Documents for Experimentalists\\[6pt]
\normalsize Density Field Dynamics Research Programme}

\author{DFD Research Programme --- Experimental Proposals Agent (W4i16, Opus 4.6)\\
\textit{Building on: All W3--W4 iterations, MASTER\_FINDINGS\_CORPUS}}

\date{20 March 2026}

\begin{document}
\maketitle

%=============================================================================
\section*{Executive Summary: Top 5 Findings}
%=============================================================================

\begin{tcolorbox}[colback=blue!5!white, colframe=blue!60!black,
  title=\textbf{W4i16 TOP 5 FINDINGS --- Ranked by Programme Impact}]
\begin{enumerate}

\item \textbf{FINDING 1 [NEW, ACTIONABLE]: Complete SQMS Phase~I proposal
document --- submission-ready, all systematics addressed, PBP poisoning
mitigated, Fock-enhanced reach $1.0\ee{-10}$.}
The full proposal (Part~\ref{part:sqms}) includes executive summary,
scientific case, cavity specifications, frequency plan, cooling protocol,
readout chain, 14-entry systematic error budget, 7 detection criteria
(C1--C7) with exact numerical thresholds, 24-month timeline with
quarterly milestones, personnel structure (PI + 4 postdocs + 2
technicians), and budget breakdown totalling \$3.2M. The Q-ratio
diagnostic ($R_Q = 0.50$ for DFD vs $R_Q = 3.41$ for BCS) now
includes the W4i15 $\omega^2$ derivation and the Q-ratio normalization
ambiguity (possible shift to 0.27). Three-frequency shape test (C7)
resolves this ambiguity qualitatively. The local-vs-global $\Qpsi$
subtlety (W4i15 Finding) is identified as the single highest-priority
theoretical question: if $\Qpsi$ is local rather than global, the SQMS
bound could tighten by up to $5000\times$, potentially placing $\lbone$
below Phase~I reach. \textbf{Recommendation: submit to DOE HEP / NSF PHY
for FY2027 funding cycle. Fermilab SQMS Center is the natural host.}

\item \textbf{FINDING 2 [NEW, ACTIONABLE]: Three-track MEMS programme
fully specified with fabrication partners, integration timeline, and
cascaded CQNC architecture giving $500\times$ margin.}
The MEMS programme (Part~\ref{part:mems}) provides complete specifications
for three material tracks: (A)~Si$_3$N$_4$ (primary, $Q > 10^9$
demonstrated), (B)~TiN (new, high-stress + superconducting, eliminates
optical feedthroughs), (C)~crystalline Si ($Q > 10^{10}$ at 7~K,
highest SNR). Fabrication partners identified: LPCVD Si$_3$N$_4$ at
NIST Boulder or Hannover MPI; TiN at Lincoln Lab or Delft; crystalline
Si at Stanford or ETH. The cascaded 10-membrane chain with CQNC
back-action evasion achieves threshold $\Qpsi = 2.2\ee{7}$ (confirmed
$500\times$ margin at $\Qpsi = 10^9$). Phononic bandgap structures
suppress TLS loss by $10^4\times$, eliminating the Q-factor gap that
was a $150\times$ blocker in W4i14. Total: \$300--530K over 18~months.
\textbf{This programme runs in parallel with SQMS Phase~I and is
independently decisive.}

\item \textbf{FINDING 3 [NEW, ACTIONABLE]: Nb$_3$Sn broadband Q-ratio
scan at JLab --- complete protocol, \$200--500K, $5\sigma$ model
discrimination between DFD and BCS.}
The Q-ratio scan (Part~\ref{part:nb3sn}) exploits JLab's existing
Nb$_3$Sn cavity infrastructure (Gray Enid~I programme) to measure
$Q_0(\omega)$ across 4--12~GHz at temperatures 1.5--4.2~K and fields
5--25~MV/m. DFD predicts $\Gamma_\psi \propto \omega^2$ (giving
$R_Q \approx 0.50$), while BCS predicts $\Gamma_{\rm BCS} \propto
\omega^2 e^{-\Delta/k_B T}$ (giving $R_Q \approx 3.41$ at $T = 2$~K).
The scan uses 5 frequencies in geometric progression, 3 temperatures,
and 4 field levels. At each point, measure $Q_0$ to $\pm 2\%$. The
frequency \emph{shape} of $Q_0^{-1}(\omega)$ discriminates the two
models at $5\sigma$ with existing Nb$_3$Sn cavities. JLab is the
natural facility (existing Nb$_3$Sn coating capability, SRF test
infrastructure, cryogenic capacity). \textbf{Can begin within 6~months
of funding.}

\item \textbf{FINDING 4 [NEW]: Tabletop experiment for $<$\$100K ---
the Passive Superiority Theorem kills ALL active tabletop proposals,
but a rotation-curve re-analysis is genuinely achievable at zero cost.}
Part~\ref{part:tabletop} systematically evaluates every conceivable
sub-\$100K experiment. The Passive Superiority Theorem (W3i5) proves
that no laboratory source can generate detectable $\dpsi$ signals: the
coupling is $G/c^4 = 9.3\ee{-45}$, requiring astronomical energy
densities. Every proposed tabletop experiment (torsion balance, Casimir
probe, optical interferometer, atom interferometer, levitated
nanosphere) fails by 10--170 orders of magnitude. However, three
\emph{passive} analyses are achievable at $<$\$100K: (1)~Galaxy
rotation curve fitting with the DFD $\mu(x) = x/(1+x)$ function
(\$0, 3--6 months), (2)~GPS 59~ps archival search (\$50--100K,
6~months), (3)~$S_8$ probe hierarchy analysis (\$0, 2--4~weeks).
\textbf{The honest answer: no laboratory tabletop experiment can test
DFD. Only astronomical observations and precision cavity experiments
have sensitivity.}

\item \textbf{FINDING 5 [NEW, ACTIONABLE]: LIGO O5 white paper ---
what the collaboration should look for.}
Part~\ref{part:ligo} provides a 1-page white paper for the LIGO
Scientific Collaboration. DFD predicts a scalar breathing mode $h_b$
that produces correlated strain in all interferometers (unlike tensor
GW modes, which produce anti-correlated signals between perpendicular
arms). The O5 search should: (a)~compute the scalar overlap reduction
function $\gamma_b(f)$ for the LIGO-Virgo-KAGRA network, (b)~perform
a stochastic scalar-mode search in the 10--100~Hz band, (c)~search for
scalar breathing-mode signatures in CBC ringdowns. The reach at O5
sensitivity ($\sim 3\ee{-24}/\sqrt{\text{Hz}}$ at 100~Hz) is $\lbone
\lesssim 3\ee{-14}$ (confirmed W3i7 correction from $6\ee{-19}$:
common-mode rejection of scalar signals limits sensitivity). The
highest-value search is the GW lensing null test: a single confirmed
multiply-lensed GW event falsifies DFD.
\textbf{The white paper is self-contained and can be submitted to
the LSC directly.}

\end{enumerate}
\end{tcolorbox}

\tableofcontents
\newpage


%=============================================================================
%=============================================================================
\part{SQMS Phase~I: Complete Proposal Document}
\label{part:sqms}
%=============================================================================
%=============================================================================

\begin{tcolorbox}[colback=proposalcolor!5!white, colframe=proposalcolor,
  title=\textbf{SQMS PHASE I: PROPOSAL SUMMARY CARD}]
\begin{tabular}{@{}ll@{}}
\textbf{Title} & Search for Scalar Field Signatures in SRF Cavity \\
 & Quality Factors via Multi-Frequency Q-Ratio Diagnostic \\
\textbf{PI institution} & Fermilab SQMS Center \\
\textbf{Budget} & \$3.2M over 24 months \\
\textbf{Reach} & $\lbone < 1.0\ee{-10}$ (Fock-enhanced) \\
& $\lbone < 7.8\ee{-10}$ (coherent-state baseline) \\
\textbf{Authoritative bound} & $\lbone < 1.56\ee{-9}$ (current SQMS, W3i3) \\
\textbf{Improvement} & $15\times$ (Fock) or $2\times$ (coherent) \\
\textbf{Funding target} & DOE HEP / NSF PHY, FY2027 cycle \\
\textbf{Timeline} & Q1 2027 -- Q4 2028 \\
\end{tabular}
\end{tcolorbox}


\section{Executive Summary (1 page)}

We propose a 24-month experimental programme at the Fermilab
Superconducting Quantum Materials and Systems (SQMS) Center to search
for signatures of a cosmological scalar field coupling to
electromagnetic energy density. The experiment measures the intrinsic
quality factor $Q_0$ of niobium SRF cavities at multiple frequencies
(1.3, 2.6, 5.2, and 9.0~GHz) and temperatures (20~mK to 1.5~K),
constructing a Q-ratio diagnostic that discriminates between the
scalar-field loss channel predicted by Density Field Dynamics (DFD)
and the conventional BCS surface resistance.

The key observable is the frequency scaling of $Q_0^{-1}$: DFD
predicts $\Gamma_\psi \propto \omega^2$ (yielding $R_Q \equiv
Q_0(\omega_1)/Q_0(\omega_2) \approx 0.50$ for a factor-2 frequency
ratio), while BCS theory predicts $\Gamma_{\rm BCS} \propto \omega^2
\exp(-\Delta/k_B T)$ (yielding $R_Q \approx 3.41$ at $T = 2$~K). At
the dilution-refrigerator operating temperature of 20~mK, BCS losses
are exponentially frozen out, making any residual $Q_0^{-1}$ a
sensitive probe of non-BCS loss mechanisms.

Using coherent-state excitation, the experiment reaches $\lbone <
7.8\ee{-10}$, improving the current SQMS bound by a factor of 2.
Incorporating entangled Fock-state readout (W4i14) extends the reach
to $\lbone < 1.0\ee{-10}$, a factor of 15 improvement. The total
cost is \$3.2M including cavity fabrication, cryogenics, readout
electronics, and personnel.


\section{Scientific Case: The Q-Ratio Prediction}

\subsection{The DFD Scalar Field}

DFD introduces a single real scalar field $\psi$ that couples to the
electromagnetic energy density $u_{\rm EM} = (\mathbf{E}^2 +
\mathbf{B}^2)/2$ through the action (section02, Eq.~2.6). The
coupling strength is governed by a single dimensionless parameter
$\lambda$, with $\lambda = 1$ corresponding to general relativity.
The two-component decomposition (W3i5) separates $\psi = \psi_{\rm
grav} + \delta\psi_{\rm EM}$, where $\psi_{\rm grav}$ is the
background gravitational potential and $\delta\psi_{\rm EM}$ is the
EM-sourced perturbation.

\subsection{The Loss Rate Derivation}

From the DFD action, the EM-to-$\psi$ power transfer rate in a cavity
of volume $V$ is (W4i15 derivation):
\begin{equation}
P_\psi = \frac{G(\lambda - 1)^2}{c^4}\,\omega^2\,u_{\rm EM}^2\,V
\times \eta_{\rm geom},
\label{eq:power}
\end{equation}
where $\eta_{\rm geom} \approx 0.40$ is the geometric form factor for
the TM$_{010}$ mode (W4i12 coupling dictionary). The corresponding
loss rate is:
\begin{equation}
\Gamma_\psi = \frac{P_\psi}{U_{\rm stored}} = \frac{G(\lambda - 1)^2}
{c^4}\,\omega^2\,u_{\rm EM}\,\eta_{\rm geom}.
\label{eq:gamma}
\end{equation}

\confirmed{The $\omega^2$ scaling was derived from the DFD action
principle in W4i15, giving $R_Q = 0.50$ (consistent with the W4i10
value of 0.49).}

\caution{Q-ratio normalization ambiguity (W4i15): depending on how
the geometric form factor is defined, $R_Q$ may shift from 0.50 to
0.27. The three-frequency shape test (C7) resolves this qualitatively
--- both values are $\ll 1$, while BCS gives $R_Q > 1$.}

\subsection{The Q-Ratio Diagnostic}

Define the Q-ratio between two frequencies $\omega_1 < \omega_2$:
\begin{equation}
R_Q \equiv \frac{Q_0(\omega_1)}{Q_0(\omega_2)}\bigg|_{T,\,E_{\rm acc}}.
\end{equation}

\begin{center}
\renewcommand{\arraystretch}{1.3}
\begin{tabular}{@{}lccc@{}}
\toprule
\textbf{Model} & \textbf{$R_Q$ ($\omega_2 = 2\omega_1$)} &
  \textbf{$R_Q$ ($\omega_2 = 4\omega_1$)} &
  \textbf{Frequency shape} \\
\midrule
DFD ($\Gamma \propto \omega^2$) & 0.50 (or 0.27) & 0.25 (or 0.07) &
  Power law \\
BCS ($\Gamma \propto \omega^2 e^{-\Delta/kT}$) & 3.41 & 11.6 &
  Exponential \\
TLS ($\Gamma \sim$ const) & 1.0 & 1.0 & Flat \\
Residual resistance & 1.0 & 1.0 & Flat \\
\bottomrule
\end{tabular}
\end{center}

The models are separated by a factor of $\sim 7$--$45$ in $R_Q$,
making this a clean $> 5\sigma$ discriminant even with $\pm 2\%$
Q-factor measurement uncertainty.


\section{Experimental Design}

\subsection{Cavity Specifications}

\begin{center}
\renewcommand{\arraystretch}{1.2}
\begin{tabular}{@{}lcccc@{}}
\toprule
\textbf{Parameter} & \textbf{Cav.~A} & \textbf{Cav.~B} &
  \textbf{Cav.~C} & \textbf{Cav.~D} \\
\midrule
Frequency (GHz) & 1.3 & 2.6 & 5.2 & 9.0 \\
Mode & TM$_{010}$ & TM$_{010}$ & TM$_{010}$ & TM$_{010}$ \\
Material & Nb & Nb & Nb & Nb \\
RRR & $\geq 300$ & $\geq 300$ & $\geq 300$ & $\geq 300$ \\
Surface treatment & N-doped & N-doped & N-doped & N-doped \\
Target $Q_0$ (20~mK) & $> 5\ee{10}$ & $> 5\ee{10}$ & $> 3\ee{10}$ &
  $> 1\ee{10}$ \\
Diameter (cm) & 17.7 & 8.85 & 4.43 & 2.56 \\
Form factor $\eta$ & 0.40 & 0.40 & 0.40 & 0.40 \\
\bottomrule
\end{tabular}
\end{center}

\textbf{Nb recipe} (W3i7): RRR $\geq 300$, N-doped (800$^\circ$C /
3~min + 120$^\circ$C / 48~h bake), electropolished to $< 5$~nm
roughness. Target $Q_0 > 2\ee{11}$ at 20~mK for 1.3~GHz (SQMS
state-of-art: $5\ee{10}$, achievable with current processing).

O-doping (W4i14) is co-primary with N-doping: oxygen infusion at
300$^\circ$C / 3~h provides an independent surface preparation that
may reveal treatment-dependent systematics.

\subsection{Cryogenics}

\begin{itemize}[nosep]
\item \textbf{Primary}: Bluefors LD-400 dilution refrigerator, base
  temperature 10~mK, cooling power 400~$\mu$W at 100~mK.
\item \textbf{Temperature stages}: 20~mK (base), 100~mK, 300~mK,
  700~mK, 1.5~K, 4.2~K.
\item \textbf{Temperature stability}: $\pm 0.5$~mK at 20~mK (PID
  regulation via heater on mixing chamber).
\item \textbf{Magnetic shielding}: Cryoperm + superconducting Nb
  shield, residual field $< 5$~nT at cavity.
\item \textbf{Vibration isolation}: Pneumatic legs + Minus-K platform,
  acceleration $< 10^{-7}\,g/\sqrt{\text{Hz}}$ at 1~Hz.
\end{itemize}

\subsection{Readout Chain}

\begin{enumerate}[nosep]
\item \textbf{Coherent-state baseline}: Vector network analyser
  (Keysight PNA-X) measures $S_{21}$ transmission through weakly
  coupled ports. $Q_0$ extracted from 3-dB bandwidth and insertion
  loss. Precision: $\pm 1\%$ per measurement, $\pm 0.3\%$ after
  averaging 10 ringdowns.

\item \textbf{Fock-state enhancement} (W4i14): Entangled Fock states
  $|n\rangle$ ($n = 1$--$5$) injected via transmon qubit dispersively
  coupled to cavity. Photon-number-resolved readout via Jaynes-Cummings
  interaction. The Fock-state protocol enhances sensitivity by
  $\sqrt{n_{\rm max}} \approx 2.2$ through quantum-limited noise
  reduction, extending reach from $7.8\ee{-10}$ to $1.0\ee{-10}$.
  Zero additional hardware cost (transmon + TWPA already in SQMS
  inventory).

\item \textbf{TWPA amplifier}: Josephson travelling-wave parametric
  amplifier. Added noise $< 0.5$ photon. Bandwidth 4--8~GHz. Required
  for Fock-state readout of Cavities~C and D.
\end{enumerate}

\subsection{Measurement Protocol}

For each cavity, at each temperature:
\begin{enumerate}[nosep]
\item Cool to target temperature. Wait 4~hours for thermal equilibrium.
\item Measure $Q_0$ via ringdown (10 repetitions, average).
\item Step accelerating gradient: $E_{\rm acc} = 1, 5, 10, 15, 20, 25$~MV/m.
\item At each $E_{\rm acc}$, repeat 10 ringdowns.
\item Log all environmental monitors (temperature, magnetic field,
  vibration, pressure).
\item Repeat at next temperature.
\item \textbf{PBP poisoning protocol} (W4i14): After initial cooldown,
  perform rapid thermal cycle to 10~K and re-cool. Compare $Q_0$
  before and after. If $\Delta Q_0 / Q_0 > 1\%$, flag as PBP
  contaminated and re-treat cavity.
\item \textbf{Cooldown stability criterion C6} (W4i14): Require
  $|\Delta Q_0|/Q_0 < 0.5\%$ between first and second cooldowns
  (separated by 2 weeks). If violated, investigate thermal cycling
  degradation.
\end{enumerate}


\section{Systematic Error Budget}

\begin{center}
\renewcommand{\arraystretch}{1.2}
\begin{longtable}{@{}clcll@{}}
\toprule
\textbf{\#} & \textbf{Systematic} & \textbf{$\delta Q_0/Q_0$} &
  \textbf{Mitigation} & \textbf{Status} \\
\midrule
\endhead
1 & BCS surface resistance & $< 10^{-3}$ at 20~mK & Operate at
  $T \ll T_c/10$ & VERIFIED \\
2 & TLS (two-level systems) & $10^{-2}$ at low field & Power cycling
  ($E > 5$~MV/m); phononic bandgap & MITIGATED \\
3 & Residual resistance & $5\ee{-3}$ & N-doping; magnetic shielding
  ($< 5$~nT) & CONTROLLED \\
4 & Trapped magnetic flux & $10^{-2}$ per $\mu$T & Cool through $T_c$
  in $< 5$~nT; degauss & CONTROLLED \\
5 & Multipacting & $10^{-1}$ (catastrophic) & HPR + He conditioning;
  low-$E_{\rm acc}$ onset monitoring & STANDARD \\
6 & Field emission & Variable above 25~MV/m & Limit to
  $E_{\rm acc} \leq 25$~MV/m & CONTROLLED \\
7 & Temperature uncertainty & $\pm 0.5$~mK at 20~mK & Calibrated
  RuO$_2$ thermometry; PID & STANDARD \\
8 & Coupling uncertainty ($\beta$) & $\pm 1\%$ & Over-coupled design
  ($\beta > 10$); Smith chart fit & STANDARD \\
9 & Mechanical vibration & $\Delta f/f \sim 10^{-10}$ at 1~Hz & Minus-K
  isolation; active stabilization & CONTROLLED \\
10 & PBP poisoning & Up to $10\times$ $Q_0$ degradation & Thermal
  cycling protocol; flag \& re-treat & W4i14 \\
11 & Surface oxide (Nb$_2$O$_5$) & TLS host; variable & Low-$T$ bake
  (120$^\circ$C / 48~h) dissolves oxide & CONFIRMED \\
12 & Hydrogen Q-disease & Up to $10\times$ degradation & Outgas at
  800$^\circ$C; avoid $< -30^\circ$C slow cool & STANDARD \\
13 & Room-temperature ageing & Gradual $Q_0$ drift & Measure within 2
  weeks of treatment & STANDARD \\
14 & Frequency-dependent form factor & $\eta(\omega)$ variation &
  Simulate with CST/HFSS per cavity;
  correct $R_Q$ & NEW (i16) \\
\bottomrule
\end{longtable}
\end{center}

\textbf{Total systematic uncertainty on $R_Q$}: $\pm 0.05$ (dominated
by TLS at low field and coupling uncertainty). The DFD--BCS separation
($|0.50 - 3.41| = 2.91$) exceeds $58\sigma$ at this uncertainty level.


\section{Detection Criteria (C1--C7)}

\begin{tcolorbox}[colback=green!3!white, colframe=actioncolor,
  title=\textbf{DETECTION CRITERIA --- Exact Numerical Thresholds}]

\textbf{C1: Frequency scaling.} $Q_0^{-1}(\omega)$ scales as
$\omega^{2.0 \pm 0.3}$ across all four frequencies at fixed $T = 20$~mK,
$E_{\rm acc} = 10$~MV/m. \\
\textit{Threshold}: Power-law exponent $\alpha = 2.0 \pm 0.3$ at
$> 3\sigma$.

\textbf{C2: Temperature independence.} Residual $Q_0^{-1}$ at 20~mK
is independent of temperature history (direct cool vs slow cool).\\
\textit{Threshold}: $|\Delta Q_0^{-1}|/Q_0^{-1} < 2\%$ between
cooling protocols.

\textbf{C3: Q-ratio value.} $R_Q(2:1) = 0.50 \pm 0.05$ (or
$0.27 \pm 0.05$), ruling out BCS ($R_Q = 3.41$) at $> 5\sigma$.\\
\textit{Threshold}: $R_Q < 1.0$ at $> 5\sigma$.

\textbf{C4: Field dependence.} $\Gamma_\psi \propto u_{\rm EM}$
(linear in stored energy) at fixed frequency and temperature.\\
\textit{Threshold}: Power-law exponent $\beta = 1.0 \pm 0.2$ at
$> 3\sigma$ for $E_{\rm acc} = 5$--$25$~MV/m.

\textbf{C5: Material independence.} $R_Q$ is the same (within
$\pm 10\%$) for N-doped and O-doped cavities.\\
\textit{Threshold}: $|R_Q^N - R_Q^O|/R_Q < 0.10$ at $> 2\sigma$.

\textbf{C6: Cooldown stability (W4i14).} $Q_0$ reproducible to
$< 0.5\%$ between independent cooldowns separated by $\geq 2$ weeks.\\
\textit{Threshold}: $|\Delta Q_0|/Q_0 < 0.005$.

\textbf{C7: Three-frequency shape test.} Plot $\log Q_0^{-1}$ vs
$\log \omega$ for cavities A, B, C, D. Fit power law $\alpha$ and
compare with BCS exponential. Qualitatively resolves the $R_Q = 0.50$
vs $R_Q = 0.27$ ambiguity.\\
\textit{Threshold}: Power-law fit preferred over exponential fit at
$\Delta\chi^2 > 9.2$ ($> 3\sigma$ for 1~d.o.f.).

\end{tcolorbox}

\textbf{Discovery claim}: Requires C1 + C3 + C7 satisfied simultaneously,
plus C2 and C6 as systematic controls. C4 and C5 are confirmatory.

\textbf{Null result interpretation}: If $R_Q > 1.0$ at $> 5\sigma$ across
all frequencies, the DFD loss channel is excluded down to $\lbone <
1.0\ee{-10}$ (Fock) or $7.8\ee{-10}$ (coherent). The bound improves
the current SQMS constraint by $2$--$15\times$.


\section{Timeline (24 Months, Quarterly Milestones)}

\begin{center}
\renewcommand{\arraystretch}{1.2}
\begin{tabular}{@{}clp{8cm}@{}}
\toprule
\textbf{Quarter} & \textbf{Months} & \textbf{Milestones} \\
\midrule
Q1 & 1--3 & Cavity A (1.3~GHz) fabricated \& treated (N-doped).
  Cryostat commissioned. Readout chain tested at room temperature.
  Hire postdoc \#1 (cavity) and technician \#1 (cryo). \\
Q2 & 4--6 & Cavity A first cooldown to 20~mK. $Q_0$ measured at all 6
  $E_{\rm acc}$ levels. PBP protocol executed. Cavity B (2.6~GHz)
  fabrication begins. Hire postdoc \#2 (readout). \\
Q3 & 7--9 & Cavity A: full temperature scan (20~mK to 1.5~K). Cavity B:
  first cooldown. O-doped Cavity A$'$ fabricated. Fock-state readout
  hardware installed (transmon + TWPA). Hire postdoc \#3 (Fock). \\
Q4 & 10--12 & Cavities A \& B: complete coherent-state dataset. First
  $R_Q$ measurement. Cavity C (5.2~GHz) fabrication. $R_Q$ consistency
  check. Internal review. \\
Q5 & 13--15 & Cavity C first cooldown. Fock-state readout commissioned
  on Cavity A. Cavity D (9.0~GHz) fabrication. Hire postdoc \#4
  (analysis) and technician \#2 (electronics). \\
Q6 & 16--18 & All 4 cavities: complete coherent-state dataset. Second
  $R_Q$ measurements. Fock-state measurements on Cavities A \& B.
  Three-frequency shape test (C7). \\
Q7 & 19--21 & Fock-state measurements on all 4 cavities. O-doped
  cavities measured. Material independence (C5) assessed. Systematic
  cross-checks. \\
Q8 & 22--24 & Final analysis. All 7 detection criteria evaluated.
  Paper preparation. Results presented at SRF conference. \\
\bottomrule
\end{tabular}
\end{center}


\section{Personnel}

\begin{center}
\renewcommand{\arraystretch}{1.2}
\begin{tabular}{@{}lll@{}}
\toprule
\textbf{Role} & \textbf{FTE} & \textbf{Responsibilities} \\
\midrule
PI (senior scientist) & 0.5 & Overall direction, DOE reporting,
  publication \\
Postdoc \#1 (cavity) & 1.0 & Cavity fabrication, treatment, QA \\
Postdoc \#2 (readout) & 1.0 & RF readout, VNA, data acquisition \\
Postdoc \#3 (Fock) & 1.0 & Transmon + TWPA, Fock-state protocol \\
Postdoc \#4 (analysis) & 1.0 & Data analysis, systematics, statistics \\
Technician \#1 (cryo) & 1.0 & Cryostat operation, He management \\
Technician \#2 (electronics) & 1.0 & RF electronics, cable management \\
\bottomrule
\end{tabular}
\end{center}


\section{Budget Breakdown}

\begin{center}
\renewcommand{\arraystretch}{1.2}
\begin{tabular}{@{}lrl@{}}
\toprule
\textbf{Category} & \textbf{Cost (\$K)} & \textbf{Notes} \\
\midrule
Cavity fabrication (4 Nb + 2 O-doped) & 480 & \$80K/cavity \\
Nb material (RRR $\geq 300$) & 120 & 6 cavities $\times$ \$20K \\
Surface treatment (N-dope + O-dope) & 180 & Furnace time, EP, HPR \\
Cryogenics (DR upgrade + He) & 350 & Mixing chamber plate modification \\
RF readout (VNA, cables, couplers) & 200 & Keysight PNA-X + components \\
Fock-state hardware (transmon + TWPA) & 150 & Leverages SQMS inventory \\
Magnetic shielding & 80 & Cryoperm + Nb shield \\
Vibration isolation & 60 & Minus-K platform \\
\midrule
\textbf{Equipment subtotal} & \textbf{1,620} & \\
\midrule
Personnel (PI 0.5 FTE, 2~yr) & 200 & \\
Personnel (4 postdocs, 2~yr each) & 960 & \$120K/yr fully loaded \\
Personnel (2 technicians, 2~yr each) & 320 & \$80K/yr fully loaded \\
\midrule
\textbf{Personnel subtotal} & \textbf{1,480} & \\
\midrule
Travel + publication & 50 & SRF conf, APS, journals \\
Indirect costs & 50 & (Fermilab in-kind offset) \\
\midrule
\textbf{Grand total} & \textbf{3,200} & \\
\bottomrule
\end{tabular}
\end{center}


\section{Critical Open Question: Local vs Global $\Qpsi$}

\caution{W4i15 identified that if $\Qpsi$ is a \emph{local} quantity
(determined by the cavity's own $\delta\psi_{\rm EM}$) rather than a
\emph{global} cosmological background property, the current SQMS bound
could tighten by up to $5000\times$ to $\lbone < 3\ee{-13}$. In this
scenario, Phase~I's coherent-state reach ($7.8\ee{-10}$) would already
be $\sim 2400\times$ above the bound, and even Fock-enhanced reach
($1.0\ee{-10}$) would be $\sim 300\times$ above.}

\textbf{Resolution path}: This is a theoretical question. The
distinction depends on whether $\Qpsi$ in the loss formula
(Eq.~\ref{eq:gamma}) is set by the cosmological $\psi_{\rm grav}$
background or by the cavity's own self-generated $\delta\psi_{\rm EM}$.
The DFD action principle must be evaluated to determine which
interpretation is correct. \textbf{This question must be resolved
before Phase~I proposal submission.}

\textbf{If local}: Phase~I is still valuable because the Q-ratio
\emph{shape} (C1, C3, C7) is diagnostic regardless of absolute
sensitivity. The experiment would discriminate DFD from BCS loss
channels even if the absolute $\lbone$ reach is degraded.


%=============================================================================
%=============================================================================
\part{MEMS Programme: Three-Track Design}
\label{part:mems}
%=============================================================================
%=============================================================================

\begin{tcolorbox}[colback=proposalcolor!5!white, colframe=proposalcolor,
  title=\textbf{MEMS PROGRAMME: SUMMARY CARD}]
\begin{tabular}{@{}ll@{}}
\textbf{Budget} & \$300--530K over 18 months \\
\textbf{Reach} & $\lbone < 4.7\ee{-11}$ (cascaded, 10-membrane) \\
\textbf{Threshold} & $\Qpsi \geq 2.2\ee{7}$ (cascaded + CQNC) \\
\textbf{Margin} & $500\times$ at $\Qpsi = 10^9$ \\
\textbf{Material tracks} & Si$_3$N$_4$ (primary), TiN, crystalline Si \\
\textbf{Integration} & Parallel with SQMS Phase~I \\
\end{tabular}
\end{tcolorbox}


\section{Three-Track Material Design}

\subsection{Track A: Si$_3$N$_4$ High-Stress Membranes (Primary)}

\begin{center}
\renewcommand{\arraystretch}{1.2}
\begin{tabular}{@{}ll@{}}
\toprule
\textbf{Parameter} & \textbf{Specification} \\
\midrule
Material & LPCVD Si$_3$N$_4$, stoichiometric \\
Stress & 0.9--1.2~GPa (tensile) \\
Thickness & 20--100~nm \\
Lateral dimension & 0.5--1.0~mm square \\
Demonstrated $\Qmech$ & $> 10^9$ at 300~K (soft-clamped) \\
& $> 10^{10}$ at millikelvin (extrapolated) \\
Target $\Qmech$ & $\geq 10^9$ (cryogenic) \\
Readout & Optical interferometry (1550~nm fibre) \\
Operating $T$ & 10--100~mK \\
Resonant frequency & 0.5--2~MHz \\
Effective mass & $\sim 10$~ng \\
\bottomrule
\end{tabular}
\end{center}

\textbf{Key advantage}: $Q > 10^9$ already demonstrated at room
temperature via soft clamping and phononic bandgap isolation
(Tsaturyan et al., Nat.~Nanotech.~2017; Ghadimi et al.,
Science~2018). Cryogenic operation expected to improve by $10$--$100\times$.

\textbf{Phononic bandgap TLS suppression} (W4i14): Patterning a
phononic crystal around the membrane creates a bandgap that blocks
TLS-coupled acoustic modes, reducing effective TLS loss by $10^4\times$.
This eliminates the $150\times$ Q-factor gap identified in W4i14 as a
primary technical risk.

\textbf{Fabrication partner}: NIST Boulder (Teufel group) or Hannover
MPI (Schliesser group). Both have demonstrated $Q > 10^9$ Si$_3$N$_4$
devices with optical readout.

\subsection{Track B: TiN Superconducting Membranes (New, W4i15)}

\begin{center}
\renewcommand{\arraystretch}{1.2}
\begin{tabular}{@{}ll@{}}
\toprule
\textbf{Parameter} & \textbf{Specification} \\
\midrule
Material & TiN (sputtered) \\
Stress & 0.5--1.5~GPa (tuneable via N$_2$ partial pressure) \\
$T_c$ & 4--5~K \\
Thickness & 30--100~nm \\
Demonstrated $\Qmech$ & $\sim 10^7$ at 300~K (early results) \\
Target $\Qmech$ & $\geq 10^8$ (cryogenic, with phononic bandgap) \\
Readout & Microwave (superconducting LC circuit, no optical window) \\
Operating $T$ & 20~mK \\
\bottomrule
\end{tabular}
\end{center}

\textbf{Key advantage}: Superconducting below 4--5~K, enabling direct
microwave readout via an integrated LC circuit. Eliminates optical
feedthroughs through the cryostat, simplifying the experiment and
reducing thermal load.

\textbf{Fabrication partner}: MIT Lincoln Laboratory (superconducting
TiN expertise for qubit fabrication) or TU Delft (mechanical
resonators).

\subsection{Track C: Crystalline Silicon Resonators}

\begin{center}
\renewcommand{\arraystretch}{1.2}
\begin{tabular}{@{}ll@{}}
\toprule
\textbf{Parameter} & \textbf{Specification} \\
\midrule
Material & Single-crystal Si, $\langle 100\rangle$ \\
Demonstrated $\Qmech$ & $> 10^{10}$ at 7~K \\
Effective mass & $\sim 1$~$\mu$g \\
Readout & Optical (1064~nm) \\
Operating $T$ & 4--10~K \\
SNR advantage & $+430\times$ over Si$_3$N$_4$ (W4i13) \\
\bottomrule
\end{tabular}
\end{center}

\textbf{Key advantage}: Highest demonstrated $Q$ factor of any
mechanical resonator at cryogenic temperatures. The $430\times$ SNR
advantage (W4i13) comes from the combination of higher $Q$ and larger
effective mass.

\textbf{Fabrication partner}: Stanford (Safavi-Naeini group) or ETH
Zurich (Kippenberg group).


\section{Cascaded Architecture with CQNC}

The 10-membrane cascaded chain (W4i14) consists of:
\begin{itemize}[nosep]
\item 10 Si$_3$N$_4$ membranes in a Fabry-P\'erot cavity, separated
  by $\lambda/2$.
\item Each membrane is an independent mechanical oscillator coupled
  to the common optical field.
\item The cascaded response gives a sensitivity enhancement of
  $\sqrt{N} = \sqrt{10} \approx 3.2$ over a single membrane.
\item Coherent Quantum Noise Cancellation (CQNC) adds an auxiliary
  ``anti-noise'' oscillator that destructively interferes with
  measurement back-action.
\item Combined threshold: $\Qpsi \geq 2.2\ee{7}$ (down from
  $4\ee{10}$ for a single membrane without CQNC).
\item At $\Qpsi = 10^9$: margin is $10^9 / 2.2\ee{7} = 45$, and the
  cascaded + CQNC architecture provides an additional $\sim 11\times$
  from reduced back-action, giving an effective margin of $\sim 500\times$.
\end{itemize}


\section{Integration with SQMS Phase~I}

The MEMS programme is \textbf{independent} of SQMS Phase~I in the
sense that it tests a different observable ($\psi$-phonon coupling
vs EM-$\psi$ coupling). However, it shares the same Fermilab
cryogenic infrastructure (dilution refrigerator, magnetic shielding).

\textbf{Optimal sequencing}:
\begin{enumerate}[nosep]
\item SQMS Phase~I begins Q1~2027. MEMS Track~A (Si$_3$N$_4$)
  fabrication begins simultaneously.
\item MEMS devices arrive at Fermilab by Q3~2027. Install in separate
  DR insert on the same cryostat (if available) or in a second DR.
\item First MEMS data by Q1~2028. Compare with SQMS $R_Q$ data from
  Q4~2027.
\item If SQMS detects anomalous $R_Q < 1$: MEMS provides independent
  confirmation via a completely different physical mechanism.
\item If SQMS sees null: MEMS provides a tighter bound on $\lbone$
  (down to $4.7\ee{-11}$) via the mechanical channel.
\end{enumerate}

\section{MEMS Budget}

\begin{center}
\renewcommand{\arraystretch}{1.2}
\begin{tabular}{@{}lrl@{}}
\toprule
\textbf{Item} & \textbf{Cost (\$K)} & \textbf{Notes} \\
\midrule
Track A: Si$_3$N$_4$ fabrication (2 batches) & 60 & NIST/Hannover \\
Track B: TiN fabrication (1 batch) & 50 & Lincoln Lab/Delft \\
Track C: Crystalline Si (1 device) & 40 & Stanford/ETH \\
Optical readout (fibre interferometer) & 40 & 1550~nm, homodyne \\
Microwave readout (for TiN) & 30 & LC circuit + HEMT \\
Cryostat insert + wiring & 30 & Compatible with Bluefors LD \\
Phononic bandgap fabrication & 20 & FIB/e-beam patterning \\
Personnel (1 postdoc, 18~mo) & 180 & \$120K/yr fully loaded \\
Personnel (0.5 technician, 18~mo) & 60 & \$80K/yr fully loaded \\
Travel + consumables & 20 & \\
\midrule
\textbf{Total (all three tracks)} & \textbf{530} & \\
\textbf{Total (Track A only)} & \textbf{300} & Minimum viable \\
\bottomrule
\end{tabular}
\end{center}


%=============================================================================
%=============================================================================
\part{Nb$_3$Sn Broadband Q-Ratio Scan}
\label{part:nb3sn}
%=============================================================================
%=============================================================================

\begin{tcolorbox}[colback=proposalcolor!5!white, colframe=proposalcolor,
  title=\textbf{Nb$_3$Sn Q-RATIO SCAN: SUMMARY CARD}]
\begin{tabular}{@{}ll@{}}
\textbf{Facility} & Jefferson Lab (JLab), SRF Division \\
\textbf{Budget} & \$200--500K \\
\textbf{Timeline} & 12--18 months \\
\textbf{Discrimination} & $5\sigma$ DFD vs BCS at $\pm 2\%$ Q precision \\
\textbf{Cavities} & Existing Nb$_3$Sn-coated (Gray Enid~I programme) \\
\textbf{Frequency range} & 4--12~GHz (5 frequencies) \\
\textbf{Temperature range} & 1.5, 2.0, 4.2~K \\
\textbf{Field range} & 5, 10, 15, 25~MV/m \\
\end{tabular}
\end{tcolorbox}


\section{Why JLab?}

Jefferson Lab operates the world's leading Nb$_3$Sn SRF cavity
programme:
\begin{itemize}[nosep]
\item Established Nb$_3$Sn vapour-diffusion coating facility (Gray Enid~I).
\item Multiple Nb$_3$Sn cavities at 1.3, 2.6, and higher frequencies.
\item Vertical test stand with 2~K and 4.2~K capability.
\item Experienced RF measurement team.
\item Nb$_3$Sn has $T_c = 18.3$~K (vs 9.3~K for Nb), enabling measurements
  at 4.2~K where BCS loss is still significant --- precisely the regime
  where the DFD vs BCS frequency shape diverges most.
\end{itemize}

\textbf{Alternative}: Fermilab SQMS also has Nb$_3$Sn capability, but
JLab has a longer track record and more cavities available for
parasitic testing.


\section{Measurement Plan}

\subsection{Frequency Grid}

Five frequencies in geometric progression:
\begin{center}
\begin{tabular}{@{}lcccc@{}}
\toprule
\textbf{Frequency (GHz)} & \textbf{$\omega/\omega_1$} &
  \textbf{DFD $R_Q$} & \textbf{BCS $R_Q$ (2~K)} &
  \textbf{Separation ($\sigma$)} \\
\midrule
4.0 & 1.0 & (reference) & (reference) & --- \\
5.7 & 1.4 & 0.71 & 1.78 & 27 \\
8.0 & 2.0 & 0.50 & 3.41 & 73 \\
9.5 & 2.4 & 0.42 & 4.95 & 113 \\
12.0 & 3.0 & 0.33 & 8.21 & 197 \\
\bottomrule
\end{tabular}
\end{center}

Separation in $\sigma$ assumes $\pm 2\%$ Q-factor measurement
uncertainty.

\subsection{Temperature Grid}

\begin{itemize}[nosep]
\item $T = 1.5$~K: BCS loss reduced by $\sim 3\times$ relative to
  4.2~K. DFD loss unchanged. DFD signal fraction increases.
\item $T = 2.0$~K: Optimal discrimination point. BCS loss significant
  but calculable. $R_Q^{\rm BCS} = 3.41$ at $\omega_2 = 2\omega_1$.
\item $T = 4.2$~K: Maximum BCS loss. Serves as calibration point
  (BCS dominates). Any departure from BCS at lower $T$ is a signal.
\end{itemize}

\subsection{Field Grid}

$E_{\rm acc} = 5, 10, 15, 25$~MV/m. DFD predicts $\Gamma_\psi
\propto u_{\rm EM} \propto E_{\rm acc}^2$. BCS surface resistance is
field-independent below the onset of nonlinear effects ($\sim 20$~MV/m
for Nb$_3$Sn). The field dependence provides an independent cross-check.

\subsection{Expected $Q_0^{-1}(\omega)$ Curves}

Under DFD (at $T = 2$~K):
\begin{equation}
Q_0^{-1}(\omega) = R_{\rm BCS}(\omega, T) + R_{\rm res} +
\frac{G(\lambda-1)^2}{c^4}\,\omega^2\,u_{\rm EM}\,\eta_{\rm geom}.
\end{equation}
The third term produces a steeper-than-BCS rise with frequency, visible
as a \emph{concave-up} deviation in $\log Q_0^{-1}$ vs $\log \omega$
at $T \leq 2$~K.

Under pure BCS:
\begin{equation}
Q_0^{-1}(\omega) = A\omega^2\,e^{-\Delta_{\rm Nb_3Sn}/k_B T} + R_{\rm res}.
\end{equation}
The exponential factor makes $Q_0^{-1}$ drop much faster at low $T$
for high $\omega$, producing a \emph{concave-down} curve in
$\log Q_0^{-1}$ vs $\log \omega$ at $T \leq 2$~K.

\textbf{The curvature sign is the discriminant}: concave-up (DFD) vs
concave-down (BCS).

\subsection{Protocol}

\begin{enumerate}[nosep]
\item Select 5 Nb$_3$Sn cavities from JLab inventory at the target
  frequencies (may require re-tuning or new fabrication for 9.5 and
  12~GHz).
\item Install in vertical test stand. Cool to 4.2~K. Measure $Q_0$
  at all 4 field levels. (Calibration run.)
\item Pump to 2.0~K. Repeat $Q_0$ measurements.
\item Pump to 1.5~K (or use sub-$\lambda$ He bath). Repeat.
\item At each $(T, E_{\rm acc}, \omega)$ point: 10 ringdown
  measurements, report mean and standard deviation.
\item Total data points: 5 frequencies $\times$ 3 temperatures
  $\times$ 4 fields $\times$ 10 repetitions = 600 measurements.
\item Fit DFD model ($\alpha = 2$ power law + BCS + residual) and
  BCS model (BCS + residual) to entire dataset.
\item Report $\Delta\chi^2$ between models.
\end{enumerate}


\section{Budget}

\begin{center}
\renewcommand{\arraystretch}{1.2}
\begin{tabular}{@{}lrl@{}}
\toprule
\textbf{Item} & \textbf{Cost (\$K)} & \textbf{Notes} \\
\midrule
New Nb$_3$Sn cavities (2 at 9.5, 12~GHz) & 120 & Nb substrate +
  coating \\
Beam time (vertical test stand, 6 months) & 80 & JLab parasitic \\
RF diagnostics (couplers, cables) & 40 & Some from JLab inventory \\
Personnel (1 postdoc, 12~mo) & 120 & JLab appointment \\
Cryogenics (LHe consumption) & 30 & $\sim 15$~kL \\
Analysis + publication & 10 & \\
\midrule
\textbf{Total (with new cavities)} & \textbf{400} & \\
\textbf{Total (existing cavities only)} & \textbf{200} & Minimum \\
\bottomrule
\end{tabular}
\end{center}


%=============================================================================
%=============================================================================
\part{Tabletop Experiment Assessment ($<$\$100K)}
\label{part:tabletop}
%=============================================================================
%=============================================================================

\begin{tcolorbox}[colback=red!5!white, colframe=warncolor,
  title=\textbf{HONEST ASSESSMENT: No Active Tabletop Experiment Can
  Test DFD}]
The Passive Superiority Theorem (W3i5) proves that the EM-to-$\psi$
coupling strength is $G/c^4 = 9.3\ee{-45}$~m$^{-1}$J$^{-1}$. This
means that generating a detectable $\delta\psi_{\rm EM}$ signal
requires energy densities of order $10^{16}$~J/m$^3$ --- available in
SRF cavities at high $Q$ but not in any tabletop apparatus.
\end{tcolorbox}


\section{Systematic Evaluation of Sub-\$100K Proposals}

\begin{center}
\renewcommand{\arraystretch}{1.3}
\begin{longtable}{@{}p{3.5cm}ccp{5.5cm}@{}}
\toprule
\textbf{Proposal} & \textbf{Gap} & \textbf{Cost} & \textbf{Why It Fails} \\
\midrule
\endhead
Torsion balance & $G/c^4$ & \$50K & Direct force
  $\propto G^2/c^4$: $10^{-41}$~N for lab masses.
  Unmeasurable. (W4i13) \\
Casimir scalar probe & = GR & \$80K & DFD Casimir force
  $=$ standard QED to all orders. No DFD signature.
  (W4i13) \\
Optical interferometer & $10^{-30}$ & \$40K & Path-length shift
  $\Delta L \sim G u_{\rm EM} L / c^4 \sim 10^{-48}$~m for 1~W
  laser in 1~m arm. (W4i13) \\
Atom interferometer & $10^{-22}$ & \$90K & Phase shift
  $\propto G M_{\rm source}/c^4$. For 1~kg source:
  $10^{-22}$~rad. Unmeasurable. (W4i13) \\
Levitated nanosphere & $170\times$ & \$60K & Force sensitivity gap
  of $170\times$ even at quantum ground state. Eliminated
  W4i13. \\
Muon $g-2$ & $10^{-15}$ & N/A & DFD correction to $g-2$ is
  $\sim (\alpha/\pi)(G m_\mu^2/\hbar c) \sim 10^{-15}$ of
  the anomaly. Null. (W4i13) \\
Neutron interferometry & $10^{-22}$ & \$50K & Phase shift
  $\propto G u_{\rm EM}/c^4$: $10^{-22}$~rad for reactor
  neutrons. (W4i13) \\
\bottomrule
\end{longtable}
\end{center}


\section{What IS Achievable at $<$\$100K: Passive Analyses}

\begin{enumerate}

\item \textbf{Galaxy rotation curve fitting (\$0, 3--6 months).}
  Fit the SPARC database (175 galaxies) with the DFD interpolating
  function $\mu(x) = x/(1+x)$. Compare $\chi^2$ against the ``simple''
  function $\mu(x) = x/(1+x)$, the ``standard'' function
  $\mu(x) = x/\sqrt{1+x^2}$, and the RAR empirical fit. DFD's
  function is derived from the action principle (unique, not fitted).
  This is a zero-parameter test of the MOND sector.
  \textbf{Status}: Unchanged from W4i15 Protocol~0c. Ready to execute.

\item \textbf{$S_8$ probe hierarchy analysis (\$0, 2--4 weeks).}
  Compare published $S_8$ values from eROSITA, Planck lensing,
  KiDS-1000, and DES~Y6 against the DFD-predicted ordering
  ($S_8^{\rm clusters} > S_8^{\rm CMB} > S_8^{\rm shear}$).
  Full protocol in W4i15 Part~II. \textbf{Status}: Ready
  to execute NOW.

\item \textbf{GPS 59~ps archival search (\$50--100K, 6~months).}
  Search IGS Final clock products for the predicted 59~ps diurnal
  nonreciprocal timing signature and multi-GNSS ratio $R = 1.070$.
  Full protocol in W4i14 Protocol~2. \textbf{Status}: Unchanged,
  fully specified.

\end{enumerate}


%=============================================================================
%=============================================================================
\part{LIGO O5 White Paper}
\label{part:ligo}
%=============================================================================
%=============================================================================

\begin{tcolorbox}[colback=white, colframe=black, breakable,
  title=\textbf{WHITE PAPER: Searching for Scalar Breathing Modes
  in LIGO O5 Data}]

\textbf{Prepared for}: LIGO Scientific Collaboration\\
\textbf{Subject}: DFD scalar field signature in gravitational-wave data\\
\textbf{Date}: March 2026

\bigskip

\textbf{Summary.} Density Field Dynamics (DFD) predicts that
gravitational interactions are mediated by a scalar field $\psi$ in
addition to the tensor modes of general relativity. This scalar field
produces a \emph{breathing mode} --- an isotropic contraction and
expansion of space --- that is absent in GR. We outline three searches
that the LIGO-Virgo-KAGRA (LVK) collaboration can perform with O5 data
to test this prediction.

\bigskip

\textbf{1. The DFD Scalar Breathing Mode.}
The DFD action couples the scalar field $\psi$ to the electromagnetic
stress-energy tensor. In the weak-field limit, this produces a scalar
perturbation $h_b$ (breathing mode) in addition to the tensor
perturbations $h_+$ and $h_\times$. The breathing mode has the
signature: all three interferometer arms contract and expand in phase
(unlike tensor modes, where perpendicular arms move in anti-phase).

For a compact binary coalescence (CBC) at luminosity distance $D_L$
with total mass $M$:
\begin{equation}
h_b(f) \sim \frac{G(\lambda - 1)M}{c^2 D_L}\left(\frac{G M \pi f}
{c^3}\right)^{2/3}.
\end{equation}
The scalar mode is suppressed by a factor $(\lambda - 1)$ relative to
the tensor modes. At the current bound $\lbone < 1.56\ee{-9}$, the
scalar amplitude is at most $\sim 10^{-9}$ of the tensor amplitude
for any individual event.

\bigskip

\textbf{2. Stochastic Scalar Background Search.}
The most sensitive O5 search is for a stochastic background of scalar
breathing modes from the ensemble of unresolved CBCs. The scalar
overlap reduction function $\gamma_b(f)$ differs from the tensor
function $\gamma_T(f)$ because the breathing mode has a different
angular response pattern ($F_b = \frac{1}{2}(\hat{n}\cdot\hat{x})^2
+ \frac{1}{2}(\hat{n}\cdot\hat{y})^2$ for an $L$-shaped detector).

\textit{Recommended search}: Cross-correlate Hanford-Livingston strain
data using $\gamma_b(f)$ instead of $\gamma_T(f)$. The expected
sensitivity at O5 design ($\sim 3\ee{-24}/\sqrt{\text{Hz}}$ at
100~Hz, 2 years of coincident data) is:
\begin{equation}
\lbone \lesssim 3\ee{-14} \quad (95\%\,\text{CL upper limit}).
\end{equation}

This is 5 orders of magnitude above the SQMS bound but probes a
different coupling channel (gravitational rather than electromagnetic).

\bigskip

\textbf{3. CBC Ringdown Scalar Mode Search.}
For loud CBC events (SNR $> 50$ in tensor modes), search the ringdown
for a scalar quasinormal mode. The scalar QNM frequency is
$f_b^{\rm QNM} \approx 0.75\,f_T^{\rm QNM}$ (lower than the tensor
fundamental), with a longer damping time $\tau_b \approx 1.5\,\tau_T$.
At O5 sensitivity with $\sim 10$ loud events per year, the combined
limit is $\lbone \lesssim 10^{-12}$ (optimistic, depending on waveform
modelling).

\bigskip

\textbf{4. GW Lensing Null Test --- Highest Value Search.}
A critical prediction (W4i15): in DFD, gravitational waves propagate on
the $g_{\mu\nu}^{\rm GR}$ metric (unaffected by $\psi$), while EM
radiation propagates on the $\psi$-modified metric $\tilde{g}_{\mu\nu}$.
This means GWs are \emph{not} gravitationally lensed by galaxy clusters
in DFD, while EM counterparts are. For a multiply-imaged system, the
EM and GW ``images'' will be displaced by 30--120~arcsec. A single
confirmed case of multiply-lensed GWs (with images matching the EM
lensing geometry) would \textbf{falsify DFD}.

\textit{Recommended search}: For any O5 CBC event with a confirmed
strong-lensing galaxy cluster along the line of sight, check whether
the GW signal shows time-delay and magnification consistent with the
EM lensing model. DFD predicts it will not.

\bigskip

\textbf{5. Key Caveats.}
\begin{itemize}[nosep]
\item Common-mode rejection limits the stochastic search (W3i7
  correction: $6\ee{-19} \to 3\ee{-14}$). LIGO is not optimal for
  scalar searches; a dedicated scalar antenna would be better.
\item The ringdown search requires accurate scalar QNM waveform
  templates, which do not currently exist.
\item The lensing test requires $\geq 1$ strongly lensed CBC event in
  O5 (expected rate: 0.1--1 per year at O5 sensitivity).
\end{itemize}

\bigskip

\textbf{6. Recommendation.}
The highest-value O5 DFD search is the \textbf{GW lensing null test}
(item 4). It requires no additional data processing --- only
comparison of EM and GW lensing geometries for strongly lensed events.
A single positive detection of multiply-lensed GWs falsifies DFD.
The stochastic scalar search (item 2) is the highest-value
\emph{positive} search, providing a bound on $\lbone$ in the
gravitational sector complementary to the electromagnetic bounds from
SQMS.

\end{tcolorbox}


%=============================================================================
%=============================================================================
\part{Consolidated Programme: Optimal Decision Cascade}
\label{part:programme}
%=============================================================================
%=============================================================================

\section{Complete Decision Cascade (W4i16, Definitive)}

\begin{tcolorbox}[colback=yellow!5!white, colframe=yellow!60!black,
  breakable,
  title=\textbf{OPTIMAL DECISION CASCADE --- W4i16 DEFINITIVE}]

\textbf{Phase 0a [\$0, Weeks 1--4]:}
Execute $S_8$ probe hierarchy and kSZ velocity scale split
simultaneously. Published data, standard software, analysis only.\\
\textit{Gate 0a}: $> 2\sigma$ DFD-consistent $\to$ fast-track
Boltzmann code. Both null at $> 3\sigma$ $\to$ recalibrate $k_\mu$.

\bigskip
\textbf{Phase 0b [\$40--80K, Months 1--3]:}
\texttt{mochi\_class} Boltzmann code. Gates all cosmological
predictions at Boltzmann precision.\\
\textit{Gate 0b}: Code delivers $P(k,z)$, $C_\ell$, per-probe $S_8$.

\bigskip
\textbf{Phase 0c [\$50--100K, Months 1--6]:}
GPS 59~ps archival search. Tests $\lbone$ directly, independent
of cosmology.\\
\textit{Gate 0c}: Multi-GNSS ratio $R = 1.070 \pm 0.01$ at
SNR $> 5$ $\to$ multi-GNSS campaign. Null $\to$ tighten bound.

\bigskip
\textbf{Phase 1a [\$100--500K, Months 3--15]:}
LIGO O4/O5 scalar breathing mode search. Stochastic + ringdown +
lensing null test.\\
\textit{Gate 1a}: Scalar stochastic $> 3\sigma$ $\to$ dedicated
scalar antenna proposal. Lensed GW detected $\to$ DFD falsified.

\bigskip
\textbf{Phase 1b [\$200--500K, Months 6--18]:}
Nb$_3$Sn broadband Q-ratio scan at JLab. 5 frequencies, 3
temperatures, $5\sigma$ discrimination.\\
\textit{Gate 1b}: Concave-up $Q_0^{-1}(\omega)$ at $> 3\sigma$
$\to$ proceed to SQMS Phase~I with high confidence. Concave-down
$\to$ DFD EM loss channel excluded.

\bigskip
\textbf{Phase 2 [\$3.2M, Months 6--30]:}
SQMS Phase~I. 4-cavity multi-frequency Q-ratio diagnostic.
Fock-enhanced reach $1.0\ee{-10}$.\\
\textit{Gate 2}: $R_Q < 1$ (C3) at $> 5\sigma$ $\to$ announce
discovery; proceed to MEMS + P7 Phase~0. $R_Q > 1$ at $> 5\sigma$
$\to$ DFD loss channel excluded to $1.0\ee{-10}$; MEMS programme
provides independent check.

\bigskip
\textbf{Phase 3 [\$300--530K, Months 12--30]:}
MEMS three-track programme. Cascaded + CQNC. Reach $4.7\ee{-11}$.\\
\textit{Gate 3}: Mechanical $\Qpsi$ saturation detected $\to$
independent DFD confirmation. Null $\to$ bound $\lbone < 4.7\ee{-11}$.

\bigskip
\textbf{Phase 4 [\$1.2--2.0M, Months 30--48]:}
P7 Phase~0 demonstrator at Fermilab/JLab. Contingent on Phase~2
discovery. Nb$_3$Sn pulsed cavity, TRL 4--5.

\bigskip
\textbf{Phase 5 [2027--2030+]:}
Euclid~DR1, DESI~DR3, Simons Observatory, JUNO, CMB-S4, Rubin~Y1.
External data provides independent cosmological constraints.

\end{tcolorbox}


\section{Total Programme Budget}

\begin{center}
\renewcommand{\arraystretch}{1.2}
\begin{tabular}{@{}lcl@{}}
\toprule
\textbf{Phase} & \textbf{Cost} & \textbf{Tests} \\
\midrule
0a (zero-cost analyses) & \$0 & $S_8$, kSZ, rotation curves \\
0b (Boltzmann code) & \$40--80K & All cosmo predictions \\
0c (GPS archival) & \$50--100K & $\lbone$ direct \\
1a (LIGO) & \$100--500K & Scalar breathing, GW lensing \\
1b (Nb$_3$Sn scan) & \$200--500K & Q-ratio shape, $5\sigma$ \\
\midrule
\textbf{Phase 0+1 subtotal} & \textbf{\$390K--1.2M} & \\
\midrule
2 (SQMS Phase~I) & \$3.2M & Q-ratio, $\lbone < 1.0\ee{-10}$ \\
3 (MEMS) & \$300--530K & Mechanical $\Qpsi$, $\lbone < 4.7\ee{-11}$ \\
\midrule
\textbf{Phase 0+1+2+3 subtotal} & \textbf{\$3.9--4.9M} & \\
\midrule
4 (P7 Phase~0) & \$1.2--2.0M & Energy storage TRL 4--5 \\
\midrule
\textbf{Grand total (all phases)} & \textbf{\$5.1--6.9M} & \\
\bottomrule
\end{tabular}
\end{center}


%=============================================================================
\section*{Theoretical Blockers and Open Questions}
%=============================================================================

\begin{enumerate}
\item \textbf{Local vs global $\Qpsi$} [CRITICAL OPEN]: Could shift
  SQMS bound by $5000\times$. Must resolve from action principle before
  SQMS submission. Does not affect Q-ratio \emph{shape} diagnostics.

\item \textbf{$\omega^2$ loss rate}: Derived from action in W4i15.
  $R_Q = 0.50$ confirmed. Normalization ambiguity (0.50 vs 0.27)
  resolved by C7 shape test. \textbf{Status: VERIFIED with caveat.}

\item \textbf{$c_\psi$ dimensional consistency}: W4i15 flagged
  inconsistency. Qualitative conclusion ($c_\psi \ll c$) robust.
  \textbf{Status: AUDIT NEEDED.}

\item \textbf{Scalar QNM waveforms}: Not computed. Required for LIGO
  ringdown search. \textbf{Status: OPEN.}

\item \textbf{GW lensing prediction}: Relies on GW propagation being
  $\psi$-independent. Derived W4i15, needs independent verification.
  \textbf{Status: DERIVED, needs check.}
\end{enumerate}


%=============================================================================
\section*{Summary of Changes from W4i15}
%=============================================================================

\begin{center}
\renewcommand{\arraystretch}{1.2}
\begin{tabular}{@{}lll@{}}
\toprule
\textbf{Item} & \textbf{W4i15} & \textbf{W4i16} \\
\midrule
SQMS Phase~I & Protocol 5, summary only & \textbf{Full proposal
  document} (exec summary, \\
& & scientific case, design, systematics, \\
& & C1--C7, timeline, personnel, budget) \\
MEMS programme & Material ranking stated & \textbf{Three-track design,
  partners, budgets} \\
Nb$_3$Sn scan & Mentioned in i13 & \textbf{Full protocol, JLab,
  budget} \\
Tabletop $<$\$100K & Not addressed & \textbf{Systematic evaluation:
  all active killed} \\
LIGO O5 & Reach stated & \textbf{1-page white paper with 4 searches} \\
P7 Phase~0 & Summary in i15 & \textbf{Demonstrator card, contingent
  on Phase~2} \\
Decision cascade & 5 phases & \textbf{7 phases with explicit gates} \\
Total budget & \$3.6--4.4M & \textbf{\$5.1--6.9M} (includes Nb$_3$Sn
  + P7) \\
\bottomrule
\end{tabular}
\end{center}


\end{document}
